\bookchapter{第三十六章：手機與智慧家居——讓家變聰明}

\chapterintro{手機還能幹一件越來越重要的事——當整個家的“遙控器”。從一盞燈、一個插座到一把門鎖，智慧家居正在從“極客玩具”變成“家常配置”。這一章幫你理清協議、生態、中樞和入門套件，再用一份安全清單收尾，讓你不被“智慧”反噬。}

\booksection{一、智慧家居協議：選誰很重要}

買智慧裝置前先看它支援什麼協議，這決定了它能不能被你的手機和生態“使喚”。目前主流協議有四種，各自的脾氣如下：

\begin{itemize}
  \item Matter協議：2022年由CSA（連線標準聯盟）釋出，背後是蘋果、谷歌、亞馬遜、三星等巨頭聯手。它的目標只有一個——讓所有品牌的智慧裝置說同一種語言。2023-2026年正快速成為統一標準，支援它的裝置可同時接入Apple HomeKit、Google Home、Amazon Alexa和米家，不再被單一生態綁架。
  \item Thread協議：Matter的底層網路"管道"。低功耗、自組網（某個節點斷電，訊號會自動繞路），響應快、不卡頓。買帶Thread邊界路由器的裝置（如HomePod mini、Apple TV 4K、小米中樞閘道器），體驗明顯更穩。
  \item Wi-Fi：頻寬大、能直連路由器，但功耗高、裝置多了易擠佔通道，適合相機、掃地機這類需要傳大資料的裝置。
  \item 藍芽（BLE）：配網簡單、省電，但距離短、不能自組網，適合門鎖、溫溼度計這類小資料裝置。
  \item Zigbee：老牌低功耗自組網協議，需要專用閘道器，裝置便宜穩定，但跨品牌互通差，正被Matter+Thread逐步替代。
\end{itemize}

\begin{tcbinfo}{一句話原則}
2026年買智慧裝置，優先看是否有“Matter over Thread”標識——這是當下最經得起時間檢驗的選擇。老裝置若只支援Wi-Fi或藍芽也不必急著淘汰，先用著，新購時再轉向Matter。

\end{tcbinfo}

\booksection{二、三大生態對比}

協議是“語言”，生態是“圈子”。即使都支援Matter，你日常仍要選一個主生態來搭自動化和語音控制。三大主流生態對比如下：

\begin{center}
\begin{tabularx}{\linewidth}{|X|X|X|X|}
\hline
\textbf{生態} & \textbf{優勢} & \textbf{短板} & \textbf{適用人群} \tabularnewline
\hline
Apple HomeKit & 隱私最強（本地處理為主）、與iPhone/Siri/快捷指令聯動順滑 & 裝置偏少偏貴、需要HomePod或Apple TV做中樞 & 全家蘋果裝置、重視隱私與體驗的使用者 \tabularnewline
米家（小米） & 裝置最豐富、價效比最高、小愛同學好用 & 生態相對封閉、資料多上雲、廣告偏多 & 預算有限、想一站式配齊全屋的入門使用者 \tabularnewline
華為智聯（鴻蒙） & 與鴻蒙手機/平板/手錶無縫流轉、裝置數量快速增長 & 海外裝置支援少、生態仍在追趕 & 華為手機使用者、看重"超級終端"多裝置協同 \tabularnewline
\hline
\end{tabularx}
\end{center}

\begin{tcbinfo}{選生態的決策建議}
看手機選生態：iPhone使用者優先HomeKit，華為使用者優先智聯，其他安卓使用者米家最省心。靠Matter協議，不同生態的裝置也能被同一部手機控制，不必一次性“站隊”。

\end{tcbinfo}

\booksection{三、手機當家居中樞}

在多數入門場景裡，手機就是你的“家居中樞”——透過App控制燈光、空調、窗簾、安防，還能搭建自動化場景，讓家“自己動起來”：

\begin{itemize}
  \item 裝置控制：開啟HomeKit/米家/智慧生活App，可遠端或本地開關燈、調色溫、設定空調溫度、拉合窗簾、檢視相機畫面。
  \item 自動化場景：把"條件動作"組合起來，常用的三大場景： 回家模式：定位到家+天黑開玄關燈+開空調+拉窗簾+關門鎖。
  \item 離家模式：最後一個出門（定位離家或門鎖上鎖）關燈+關空調+關插座+啟動相機+佈防。
  \item 睡眠模式：每晚23:00關全屋主燈+拉窗簾+空調切26度+開夜燈。
\end{itemize}

NFC碰一碰：在門貼一張NFC標籤，手機一碰即觸發對應場景（碰一下開門燈、再碰一下切睡眠）。iPhone用"快捷指令"自動化，安卓用"錢包"或廠商App均可寫入。

\begin{tcbpractice}{來試試：搭一個"回家模式"}
在手機智慧家居App的“自動化”裡新增一條：觸發條件選“到達某地點”（家），動作依次加“開玄關燈”“空調設26度”“窗簾關閉”。儲存後出門逛一圈，回家時看它是否自己亮燈——第一次成功的成就感，比下單新裝置還過癮。

\end{tcbpractice}

\booksection{四、入門套件推薦}

不必一次配齊，從這四件套開始，就能覆蓋八成日常需求：

\begin{itemize}
  \item 智慧音箱：小愛同學（米家生態，價效比高）/ HomePod mini（蘋果生態，兼作Thread中樞）/ Alexa（海外生態廣）。注意：音箱決定了你的語音生態，先定生態再買音箱。
  \item 智慧燈泡：Yeelight（米家/Matter，色溫亮度可調，便宜）/ Philips Hue（色彩好、穩定，但需Hue閘道器，偏貴）。注意：買前確認燈座口徑（E27/B22）和是否需要調色。
  \item 智慧插座：小米智慧插座（便宜、電量統計）/ TP-Link Kasa（海外Matter支援好）。注意：用在大功率電器（熱水器、電暖器）時務必看額定功率，別過載。
  \item 智慧門鎖：鹿客/小米（價效比、米家聯動）/ 凱迪仕（傳統鎖具大廠、機械結構紮實）。注意：確認是否支援Matter或你所在生態，留意電池續航和應急鑰匙孔位置。
\end{itemize}

\begin{tcbinfo}{選購順序建議}
先買智慧音箱（定生態） 再買智慧燈泡和插座（低門檻、立刻有反饋） 最後買門鎖和相機（涉及安全，謹慎選）。一次只加一類裝置，熟悉後再擴，避免“買回來不會用”。

\end{tcbinfo}

\booksection{五、安全注意事項}

智慧裝置越多的家，攻擊面也越大。下面四條是底線，請逐條對照：

\begin{enumerate}
  \item 改預設密碼：相機、路由器、門鎖出廠預設密碼在網上全是公開的，買到手第一件事就是改成"長且不重複"的強密碼。
  \item 及時更新韌體：廠商推送的韌體往往修復了安全漏洞，在App裡開啟"自動更新"，或至少每月手動檢查一次。
  \item 網路隔離：把IoT裝置單獨放到一個Wi-Fi（或訪客網路）上，和手機、電腦的主網路分開。萬一某個裝置被攻破，也難以橫向威脅你的重要資料。
  \item 隱私保護：相機開啟"隱私區域"遮蔽（如臥室、衛生間不拍），關閉雲端24小時錄影；不在家的相機畫面別長期開給第三方App。
\end{enumerate}

\textbf{這幾個坑別踩}\par

\begin{itemize}
  \item 別用門鎖的"遠端開鎖"功能給陌生人——優先用臨時密碼，用完即失效。
  \item 別把相機直接對著床、對著窗外——洩露的是你自己的隱私。
  \item 別圖便宜買"白牌"無品牌裝置——多半不更新韌體，等於給家裡安了個後門。
\end{itemize}

\begin{tcbsummary}{本章小結}
\begin{itemize}
  \item Matter+Thread是2023-2026年的統一方向，買新裝置優先看這個標識
  \item 三大生態各有所長：蘋果重隱私體驗、米家重價效比、華為重多裝置協同
  \item 手機可當中樞，搭好"回家/離家/睡眠"三大場景，家就能自己動起來
  \item 入門四件套：智慧音箱、燈泡、插座、門鎖，按這個順序逐步添置
  \item 改預設密碼、更新韌體、網路隔離、隱私遮蔽——四條安全底線不能省
\end{itemize}

\end{tcbsummary}
